\chapter{Discussion}\label{ch:discussion}
\section{What temporal dawn evolution adds}
The principal result is a resolved initial-value evolution connecting solar illumination, attenuation and seven-species chemistry. The reference produces large ozone changes alongside singlet-oxygen growth, with distinct vertical structures. A stationary calculation at each SZA would not retain the same chemical history. The model therefore addresses the forward temporal problem posed in the original brief, while leaving the inverse retrieval and observational comparison for future work.

The link between ozone photolysis and \singlet\ motivates interpretation of the paired fields \citep{li2020,mlynczak1993}. In the implemented equations, however, $\Delta$ also depends on direct IRA excitation, excited-state cascades and collisional loss. A decrease in ozone need not imply an immediate proportional decrease in singlet concentration. Solar production can increase as illumination improves while the accumulated ozone is being removed, and singlet oxygen can retain earlier production through its finite lifetime.

This interpretation is a consequence of the model structure, not a measured partition of the final dawn source terms. The campaign did not disable individual reactions or bands to attribute the concentration changes. It would therefore be inappropriate to label one plotted feature as uniquely caused by a particular mechanism on the basis of visual timing alone. Such causal attribution requires a separate controlled sensitivity or reaction-budget analysis.

\section{Altitude dependence and reduction validity}
The background density and collisional rates vary greatly across 50--100~km, while spherical rays illuminate different heights at different zenith angles. These model ingredients provide reasons to expect different adjustment histories rather than a single column-wide response \citep{brasseur2005}. The non-monotonic 80-km ozone case and the different seasonal ordering at 100~km are direct examples of that height dependence.

The retained mathematics also shows why an equilibrium-based reduction must be tested in its intended temporal domain. The OH partition failed because the evolving state no longer admitted a nonnegative stationary radical allocation. Peroxide elimination then failed because its dark loss denominator could vanish while production remained positive. These are two different mathematical failures, both resolved by integrating the corresponding accepted production/loss tendencies.

The addition of dynamic species should not be presented as an empirical correction to make the results look physical. It changes the reduced state space while preserving accepted chemistry. The old stationary model remains a golden reference where its algebraic solution is meaningful. The temporal model then follows physical concentrations beyond that manifold rather than insisting on a stationary allocation.

This finding has a broader methodological implication: a reduced model's algebraic domain is part of its validity, alongside its coefficients. Checking a QSSA only at a representative noon state can miss a boundary reached during dusk or dawn. The demonstrated failures concern this specified reduced network and forcing history; they do not prove that every HOx QSSA in the atmosphere must fail or that the retained fast-state QSSA is exact.

\section{Season and latitude as combined responses}
Common-SZA comparison improves interpretability by avoiding a trivial mismatch in sunrise clock time. It does not equalize night length, speed of solar-angle change, atmospheric density, temperature or initialized concentrations. Therefore seasonal and latitudinal figures describe combined responses to their specified scenarios.

The reversal of summer/March ozone ordering between 80 and 100~km directly rejects an altitude-independent enhancement description for this experiment. It does not reject a seasonal mechanism in general. Multiple source/loss terms and input prescriptions change with height, and the noon initialization carries different atomic profiles into the night. A decomposition into geometry-only, background-only and initialization-only contributions would require new experiments, which were deliberately outside the final campaign.

Likewise the equatorial excess at 90~km/SZA70 does not establish a universal equatorial ozone maximum. It is a finite-horizon result for the selected date, activity and initialization policy. The three latitude cases are illustrative model scenarios, not a statistical latitude climatology. The high-latitude SZA60 absence is a legitimate geometric limit, and comparison should respect it rather than extend a profile beyond the calculated domain.

The selected dates represent season labels but not an average over seasons. Keeping geomagnetic/solar activity quiet makes the experiments reproducible and avoids uncontrolled driver changes; it also limits historical interpretation. The background model has date/location structure, while the near-infrared incident spectrum and several reservoir profiles remain fixed prescriptions.

\section{Importance of atmospheric variability}
The identical-initial-state frozen/dynamic comparison is the strongest controlled input sensitivity in the campaign. It shows that an evolving prescribed environment materially changes the predicted dawn, with maximum relevant differences much larger than the accepted numerical discrepancies. This supports using dynamic backgrounds for the stated modelling objective rather than assuming a static reference is harmless.

The comparison does not identify whether temperature, density, collisional partners or radiation contributes most. These quantities change together within the dynamic prescription. Moreover, a time-dependent empirical background remains different from a self-consistent dynamical atmosphere. The calculation does not include continuity, vertical advection or chemical dilution, so the large sensitivity should be labelled a difference between background prescriptions within the reduced model.

Transport can affect oxygen airglow and inferred ozone; coupled photochemical--diffusive--advective studies illustrate that broader issue \citep{zhu2007}. That literature supplies a reason to investigate transport next, not a numerical correction factor to apply to the present results. No transport tendency is borrowed from a different calculation or added after the fact.

\section{Initialization as the dominant unresolved input issue}
The finite native-atom variants show that internally consistent fast initialization does not guarantee weak dependence on the chosen slow state. Upper-level ozone and singlet outputs can differ by tens of percent. The deterministic bootstrap is convenient and reproducible, but its convenience should not be mistaken for atmospheric truth.

A date/location interface must therefore expose an explicit chemical initial state or clearly label its approximate bootstrap. Automatically returning one trajectory without declaring that choice would hide an important source of uncertainty. The accepted API supports both routes and records which was used. This is a substantive scientific design decision rather than only a user-interface concern.

The multistart noon solve establishes agreement among physical convergences for the conditional fast subsystem. It does not demonstrate uniqueness of a full atmospheric chemical state, and it does not remove the dark continuum. Nor does the accepted finite horizon wash out every initialization difference. The plotted sensitivity remains a result to be carried into the thesis conclusions.

Improved initial profiles are therefore a higher priority than refining A-band spectroscopy solely for this model's present output target. The tested A1 perturbation is small downstream, while atomic initialization is material. This ordering is conditional on the reported sensitivities and intended accuracy: a future instrument-specific spectral calculation might justify a different precision hierarchy.

\section{Scope of the numerical evidence}
The BDF/Radau and interpolation comparisons support numerical convergence of the tested finite-horizon model. The output differences below approximately $0.5\%$ are useful practical evidence, but they are not confidence intervals. Tiny algebraic residuals establish that the chosen equations were evaluated consistently; they do not reduce uncertainties in reaction coefficients, atmosphere or initial state to roundoff.

The final figures are also descriptive sample summaries. Percentiles depend on saved output sampling and relevance masks. The maximum location indicates where the largest relevant saved difference occurs under that convention. It should not be recast as a universal error bound or a probability that a prediction is wrong.

No observational matchup is included. The cited OSIRIS/SABER studies establish the application context and the need for appropriate radiative/kinetic interpretation \citep{li2020,mlynczak2007}; their validation statistics are not transferred to this thesis. Implementing concentration evolution is a prerequisite for a retrieval application, not the completion of that application.

\section{Remaining limitations and a usable framework}
The omitted transport, prescribed H$_2$O/H$_2$, reference exterior ozone VMR, fixed quiet activity and approximate initialization all restrict physical interpretation. Native availability limits also matter at lower levels: missing atomic information cannot be silently converted into a validated atmospheric profile. The final model handles that convention explicitly and quantifies alternatives rather than concealing it.

Finite-horizon acceptance means that the framework successfully evolves its supplied scenario through the relevant dawn window. It does not certify a periodic daily attractor or global behaviour over arbitrarily many days. The original brief's proposed repetitive spin-up and satellite comparison remain unfulfilled broader objectives, which the final scoped conclusions acknowledge.

Further scientific review must establish the assumptions and observational tests needed for atmospheric interpretation.
